\documentclass[11pt,a4paper]{amsart}
\usepackage[T1]{fontenc}
\usepackage{lmodern}
\usepackage[margin=25mm]{geometry}
\usepackage{amsmath,amssymb,amsthm}
\usepackage[expansion=false]{microtype}
\usepackage{array,booktabs,tabularx}
\usepackage{needspace}
\setlength{\captionindent}{0pt}
\usepackage{xurl}
\usepackage[hidelinks]{hyperref}
\urlstyle{same}
\hypersetup{
  pdftitle={Erd\H{o}s Problem 928: quantitative comparisons and original formulations},
  pdfauthor={JI HO BAE},
  pdfsubject={Elementary comparison estimates and the deduction of Erd\H{o}s Problem 928 from OpenAI's joint Dickman law}
}
\newtheorem{proposition}{Proposition}
\newtheorem{corollary}[proposition]{Corollary}
\newcommand{\N}{\mathbb N}
\newcommand{\A}{\mathcal A}
\newcommand{\B}{\mathcal B}
\newcommand{\C}{\mathcal C}
\newcolumntype{L}[1]{>{\raggedright\arraybackslash}p{#1}}
\newcolumntype{Y}{>{\raggedright\arraybackslash}X}
\title[Erd\H{o}s Problem 928 and the joint Dickman law]{Erd\H{o}s Problem 928:\\
quantitative comparisons and original formulations}
\author{JI HO BAE}
\date{October 9, 2026}

\begin{document}
\begin{abstract}
We deduce the density formula $\rho(1/\alpha)\rho(1/\beta)$ for the
precise formulation of Erd\H{o}s Problem~928 from OpenAI's joint
Dickman law. Here $\rho$ is the Dickman--de Bruijn function.
For fixed $0<\alpha,\beta<1$ and real $X\ge2$, we bound the symmetric
difference between the counting sets in the problem and in the cited
theorem, restricted to $2\le n\le X$, by
$\pi(X^\alpha)+\pi((X+1)^\beta)$, where $\pi$ counts primes.
This elementary estimate accounts for both the strict inequalities and
the change in the second threshold from $n^\beta$ to $(n+1)^\beta$.
We also treat the strict formulation with both thresholds based at $n$
and compare the original statements in Erd\H{o}s's two 1976 papers and
the 1978 Erd\H{o}s--Pomerance paper. The analytic joint law is due to
OpenAI; this note provides the finite comparison estimates and the
precise correspondence among the source formulations.
\end{abstract}
\maketitle

\section{Introduction}

For an integer $n\ge2$, let $P(n)$ be the largest prime factor of $n$, and set
$P(1)=1$. Write $\rho$ for the Dickman--de Bruijn function: the continuous
function on $[0,\infty)$ determined by
\[
 \rho(u)=1\quad(0\le u\le1),\qquad
 u\rho'(u)=-\rho(u-1)\quad(u>1).
\]
The classical smooth-number limit is
\begin{equation}\label{eq:marginal}
 \lim_{X\to\infty}\frac1X
 \#\{1\le n\le X:P(n)\le X^c\}=\rho(1/c)
 \qquad(0<c<1);
\end{equation}
see de Bruijn~\cite[p.~51, (1.1)--(1.2)]{deBruijn1951}.
We also use the standard properties $\rho(u)>0$ for finite $u\ge0$
and $\rho(u)\to0$ as $u\to\infty$; see
de Bruijn~\cite[p.~52, (1.8)]{deBruijn1951} together with the defining equation.
We use $\N=\{1,2,\ldots\}$. Counting variables are integers, whereas
the cutoff $X$ is real. Natural density is the limit of the ordinary,
unweighted count up to $X$, divided by $X$.

Problem~928 in Bloom's collection~\cite{Bloom928} asks whether, for fixed
$\alpha,\beta\in(0,1)$, the integers satisfying
\begin{equation}\label{eq:problem}
 P(n)<n^\alpha,\qquad P(n+1)<(n+1)^\beta
\end{equation}
have a natural density. The accompanying discussion asks whether its
value is $\rho(1/\alpha)\rho(1/\beta)$. The page cites Erd\H{o}s's
Manitoba and Debrecen papers~\cite{Erdos1976Manitoba,Erdos1976Debrecen}
and the paper of Erd\H{o}s and Pomerance~\cite{ErdosPomerance1978}.
Section~\ref{sec:sources} compares their original statements.

Our analytic input is OpenAI's joint Dickman
law~\cite[Theorem~1.1, p.~2]{OpenAI2026}: for every fixed
$\alpha,\beta\in(0,1)$,
\begin{equation}\label{eq:JD}
 \lim_{X\to\infty}\frac1X
 \#\{2\le n\le X:P(n)\le n^\alpha,\ P(n+1)\le n^\beta\}
 =\rho(1/\alpha)\rho(1/\beta).
\end{equation}
OpenAI presents this theorem as a resolution of the
Erd\H{o}s--Pomerance joint Dickman conjecture. Section~11.3 of the
same paper also establishes the versions with thresholds
$X^\alpha,X^\beta$ and the corresponding upper-tail conditions.
The present note connects these analytic results to the precise
conventions in the cited sources. Proposition~\ref{prop:comparison}
quantifies the changes needed for~\eqref{eq:problem} and for
Erd\H{o}s's strict formulation; Section~\ref{sec:sources} explains how
the other source statements follow.

\section{Finite comparison estimates}\label{sec:comparison}

Fix $\alpha,\beta\in(0,1)$. For $X\ge2$, define
\begin{align*}
 \A_{\alpha,\beta}(X)
   &=\{n\in\N:2\le n\le X,\ P(n)<n^\alpha,
                                \ P(n+1)<(n+1)^\beta\},\\
 \B_{\alpha,\beta}(X)
   &=\{n\in\N:2\le n\le X,\ P(n)\le n^\alpha,
                                \ P(n+1)\le n^\beta\},\\
 \C_{\alpha,\beta}(X)
   &=\{n\in\N:2\le n\le X,\ P(n)<n^\alpha,
                                \ P(n+1)<n^\beta\}.
\end{align*}
The sets $\A$ and $\C$ correspond, respectively, to Problem~928 and
condition~(1) in Erd\H{o}s's Manitoba paper~\cite[p.~25]{Erdos1976Manitoba};
$\B$ is the set counted in~\eqref{eq:JD}. Let
$\pi(y)=\#\{p\le y:p\text{ is prime}\}$, and write
$S\mathbin\triangle T$ for the symmetric difference of two sets.

\begin{proposition}[Comparison of conventions]\label{prop:comparison}
For every $\alpha,\beta\in(0,1)$ and every real $X\ge2$,
\begin{align}
 \#\bigl(\A_{\alpha,\beta}(X)\mathbin\triangle
          \B_{\alpha,\beta}(X)\bigr)
 &\le \pi(X^\alpha)+\pi((X+1)^\beta),\label{eq:AB}\\
 0\le\#\B_{\alpha,\beta}(X)-\#\C_{\alpha,\beta}(X)
 &\le \pi(X^\alpha)+\pi(X^\beta).\label{eq:BC}
\end{align}
\end{proposition}

\begin{proof}
In this proof we suppress the subscripts and the argument $X$.
Suppose first that $n\in\B\setminus\A$. Since
$P(n+1)\le n^\beta<(n+1)^\beta$, the second condition defining $\A$
holds. Thus the first condition fails, giving $P(n)=n^\alpha$.
The prime $p=P(n)\le X^\alpha$ determines the only possible value
$n=p^{1/\alpha}$. Hence $\#(\B\setminus\A)\le\pi(X^\alpha)$.

Now suppose that $n\in\A\setminus\B$. The first condition defining
$\B$ holds, and failure of the second gives
\[
 n^\beta<P(n+1)<(n+1)^\beta.
\]
Because $\beta>0$, the intervals $(m^\beta,(m+1)^\beta)$ for integers
$m\ge2$ are pairwise disjoint. The map $n\mapsto P(n+1)$ is therefore
injective on $\A\setminus\B$, with values among the primes less than
$(X+1)^\beta$. Thus $\#(\A\setminus\B)\le\pi((X+1)^\beta)$.
Adding the two bounds proves~\eqref{eq:AB}.

For the remaining comparison, $\C\subseteq\B$, and every
$n\in\B\setminus\C$ satisfies
$P(n)=n^\alpha$ or $P(n+1)=n^\beta$. The first equality contributes at
most $\pi(X^\alpha)$ integers, as above. For the second equality,
$n=P(n+1)^{1/\beta}$ is determined by a prime at most $X^\beta$,
giving at most $\pi(X^\beta)$ integers. The union bound
proves~\eqref{eq:BC}.
\end{proof}

\begin{corollary}[Problem~928 and Erd\H{o}s's strict formulation]\label{cor:density}
For every fixed $\alpha,\beta\in(0,1)$,
\begin{equation}\label{eq:conclusion}
 \lim_{X\to\infty}\frac{\#\A_{\alpha,\beta}(X)}X
 =\lim_{X\to\infty}\frac{\#\C_{\alpha,\beta}(X)}X
 =\rho(1/\alpha)\rho(1/\beta).
\end{equation}
Thus the density in Problem~928 and the density $d(\alpha,\beta)$
asked for in Erd\H{o}s's Manitoba paper~\cite[p.~25]{Erdos1976Manitoba}
both exist and equal the predicted product.
\end{corollary}

\begin{proof}
For finite sets, $|\#S-\#T|\le\#(S\mathbin\triangle T)$. Since
$\pi(y)\le\lfloor y\rfloor$, both upper bounds in
Proposition~\ref{prop:comparison} are $O(X^\alpha+X^\beta)=o(X)$.
Dividing the comparison bounds by $X$ and using~\eqref{eq:JD} gives
both limits in~\eqref{eq:conclusion}.
\end{proof}

The comparison estimates are independent of the joint Dickman law:
they bound the discrepancy caused by the changes of convention for
fixed $\alpha,\beta$. The asymptotic evaluation of the count comes
from~\eqref{eq:JD}; the comparison bounds do not give a rate of
convergence in that law.

\section{Comparison with the original sources}\label{sec:sources}

We call a threshold \emph{moving} if it depends on the integer being
counted, and \emph{fixed} if it remains constant within the counting
range. Thus $n^\alpha$ and $(n+1)^\beta$ are moving thresholds, whereas
$X^\alpha$ and $X^\beta$ are fixed thresholds. The latter still vary
with the counting endpoint. Put
\[
 a(0)=1,\qquad a(t)=1-\rho(1/t)\quad(0<t\le1).
\]
Table~\ref{tab:forms} places the reference theorem beside the four
source formulations. The final column records the normalized limits
and the results from which they follow. The subsections below explain
the changes in inequalities, thresholds, and counting ranges.

\begin{table}[htbp]
\small
\caption{Reference theorem and source formulations.
In each row, $N$ counts the integers satisfying the listed conditions.
Exponents are fixed: \mbox{$0<\alpha,\beta<1$},
\mbox{$\gamma>1$}, and \mbox{$0\le t,s\le1$}.
Here $c_\gamma=\rho(\gamma)$.
Unless indicated otherwise, \mbox{$2\le n\le X$} and $X\to\infty$.}
\label{tab:forms}
\begin{tabularx}{\textwidth}{@{}L{43mm}L{45mm}Y@{}}
\toprule
Formulation and source & Conditions & Limit and justification\\
\midrule
\textbf{Reference theorem}\newline
OpenAI~\cite[Thm.~1.1, p.~2]{OpenAI2026}
 & $P(n)\le n^\alpha$\newline
   $P(n+1)\le n^\beta$
 & $\displaystyle\frac{N}{X}\to\rho(1/\alpha)\rho(1/\beta)$\newline
   Joint Dickman law~\eqref{eq:JD}\\[5pt]
\textbf{Problem~928}\newline
Bloom~\cite{Bloom928}; Erd\H{o}s's independence statement
\cite[p.~25]{Erdos1976Manitoba}
 & $P(n)<n^\alpha$\newline
   $P(n+1)<(n+1)^\beta$
 & $\displaystyle\frac{N}{X}\to\rho(1/\alpha)\rho(1/\beta)$\newline
   This note: Proposition~\ref{prop:comparison} and
   Corollary~\ref{cor:density} (set $\A$)\\[5pt]
\textbf{Erd\H{o}s's strict formulation}\newline
Manitoba paper, condition~(1)
\cite[p.~25]{Erdos1976Manitoba}
 & $P(n)<n^\alpha$\newline
   $P(n+1)<n^\beta$
 & $\displaystyle\frac{N}{X}\to\rho(1/\alpha)\rho(1/\beta)$\newline
   This note: Proposition~\ref{prop:comparison} and
   Corollary~\ref{cor:density} (set $\C$)\\[5pt]
\textbf{Fixed-threshold pair}\newline
Erd\H{o}s (Debrecen), $r=2$
\cite[pp.~272--273]{Erdos1976Debrecen}
 & $P(n+1)\le k$\newline
   $P(n+2)\le k$\newline
   $1\le n<k^\gamma$, $k\to\infty$
 & $\displaystyle\frac{N}{k^\gamma}\to\rho(\gamma)^2=c_\gamma^2$\newline
   OpenAI~\cite[(11.7), p.~80]{OpenAI2026}, with the
   change of variables in Section~\ref{subsec:debrecen}\\[5pt]
\textbf{Fixed upper tails}\newline
Erd\H{o}s--Pomerance~\cite[p.~311]{ErdosPomerance1978}
 & $P(n)\ge X^t$\newline
   $P(n+1)\ge X^s$
 & $\displaystyle\frac{N}{X}\to a(t)a(s)$\newline
   OpenAI~\cite[(11.8) and following paragraph, p.~80]{OpenAI2026},
   including equality and endpoints\\
\bottomrule
\end{tabularx}
\end{table}

\subsection{Erd\H{o}s's Manitoba paper}
On page~25 of his Manitoba paper~\cite{Erdos1976Manitoba},
Erd\H{o}s expects independence for $P(n)<n^\alpha$ and
$P(n+1)<(n+1)^\beta$. He then asks whether the integers satisfying
condition~(1), namely
$P(n)<n^\alpha$ and $P(n+1)<n^\beta$, have a density $d(\alpha,\beta)$.
The webpage therefore follows the inequalities in the independence
statement, while condition~(1) uses $n^\beta$ as its second threshold.
Corollary~\ref{cor:density} establishes the density in both formulations.

Erd\H{o}s also asks whether these strict inequalities have infinitely
many solutions for each fixed pair $\alpha,\beta>0$. The positive
value in~\eqref{eq:conclusion} gives infinitude. For parameters outside
$(0,1)$, the same conclusion follows by choosing smaller positive
parameters inside that interval.
This infinitude result was already available: as recorded on Bloom's
webpage~\cite{Bloom928}, it follows from
Schinzel~\cite[Theorem~14, p.~231]{Schinzel1967}, applied to
$(2n+1)^2-1=4n(n+1)$. The theorem gives a constant $K>0$ and infinitely
many integers $n\ge3$ such that
\[
 P(n(n+1))\le n^{K/\log\log n}.
\]
Since $K/\log\log n\to0$, these integers satisfy Erd\H{o}s's strict
inequalities for every prescribed positive pair $\alpha,\beta$, once
$n$ is sufficiently large.

\subsection{Erd\H{o}s's Debrecen paper}\label{subsec:debrecen}
To distinguish the notation from that used above, write $\gamma>1$
for the exponent denoted by $\alpha$ in Erd\H{o}s's Debrecen
paper~\cite[pp.~272--273]{Erdos1976Debrecen}. Erd\H{o}s recalls that the
number of $m\le k^\gamma$ satisfying $P(m)\le k$ is
$(c_\gamma+o(1))k^\gamma$. He then conjectures that, for each fixed integer
$r\ge1$, as $k$ tends to infinity through positive integers,
\begin{equation}\label{eq:Debrecen}
 \#\{1\le n<k^\gamma:P((n+1)\cdots(n+r))\le k\}
  =(c_\gamma^r+o(1))k^\gamma.
\end{equation}
The marginal law~\eqref{eq:marginal} identifies $c_\gamma=\rho(\gamma)$.

For the two-variable case, we use OpenAI's fixed-threshold
law~\cite[(11.7), p.~80]{OpenAI2026}: for fixed $c,d\in(0,1)$,
\begin{equation}\label{eq:fixed}
 \lim_{X\to\infty}\frac1X
 \#\{2\le m\le X:P(m)\le X^c,\ P(m+1)\le X^d\}
 =\rho(1/c)\rho(1/d).
\end{equation}
To obtain the case $r=2$ of~\eqref{eq:Debrecen}, put $X=k^\gamma$,
$m=n+1$, and $c=d=1/\gamma$. The identity
$P(m(m+1))=\max\{P(m),P(m+1)\}$ makes the prime-factor conditions
exactly those in~\eqref{eq:fixed}.

After this shift, the ranges $1\le n<X$ and $2\le m\le X$ differ
by at most one endpoint term. Its normalized contribution vanishes,
so~\eqref{eq:fixed} gives~\eqref{eq:Debrecen} for $r=2$.

Bloom's webpage~\cite{Bloom928} states a higher-order version with
$P(n+i)\le n^\alpha$ for $0\le i<r$. Its case $r=2$ is precisely
\eqref{eq:JD} with $\beta=\alpha$. The results used here concern two
consecutive integers; they do not establish the assertion for $r\ge3$.

\Needspace{10\baselineskip}
\subsection{Erd\H{o}s and Pomerance}
Erd\H{o}s and Pomerance~\cite[p.~311]{ErdosPomerance1978} formulate
independence using fixed thresholds and upper tails. For $0\le t,s\le1$,
they ask whether
\begin{equation}\label{eq:upper}
 b(t,s)=\lim_{X\to\infty}\frac1X
 \#\{2\le n\le X:P(n)\ge X^t,\ P(n+1)\ge X^s\}
\end{equation}
exists, is continuous on $[0,1]^2$, and equals $a(t)a(s)$.
Their $a(t)$ is the corresponding one-variable upper-tail limit.
The marginal law~\eqref{eq:marginal}, with the equality cases and
endpoints treated below, identifies their $a(t)$ with the function
defined above.

For $0<t,s<1$, apply inclusion--exclusion to~\eqref{eq:fixed} and the
two marginal laws. This gives the limit $a(t)a(s)$ for the strict
upper-tail conditions $P(n)>X^t$, $P(n+1)>X^s$.
Replacing $n$ by $n+1$ in a marginal count changes its counting range
by only $O(1)$ terms, so its limit is unchanged.

Allowing equality does not change the limit. For fixed $0<t\le1$,
equality $P(m)=X^t$ can occur only when $X^t$ is a prime $p$.
Then $p\mid m$, so at most
$\lfloor(X+1)/p\rfloor=O(X^{1-t}+1)=o(X)$ such $m\le X+1$ occur.

At exponent $0$ the upper-tail condition is automatic for $m\ge2$.
When just one exponent is zero, the remaining condition is evaluated
by the marginal law. At exponent $1$, the condition $P(m)\ge X$ in
the range $2\le m\le X+1$ forces $m\in[X,X+1]$, so at most two
integers occur. Thus~\eqref{eq:upper} holds at all endpoints as well.
The continuity of $\rho$, together with $\rho(1)=1$ and
$\rho(u)\to0$ as $u\to\infty$, shows that $a$ is continuous on $[0,1]$.
Hence $b(t,s)=a(t)a(s)$ is continuous on the whole square.

OpenAI already gives this deduction, including the equality cases and
endpoints, in~\cite[(11.8) and following paragraph, p.~80]{OpenAI2026}.
The purpose here is to identify its exact correspondence with the
Erd\H{o}s--Pomerance formulation~\cite{ErdosPomerance1978}. Thus all four
source formulations in Table~\ref{tab:forms} follow from the cited
joint Dickman law and its fixed-threshold version, through the
comparisons above.

\begin{thebibliography}{9}
\bibitem{Bloom928}
T.~F. Bloom, \emph{Erd\H{o}s Problem \#928},
accessed October~9, 2026.
\newline\url{https://www.erdosproblems.com/928}.

\bibitem{deBruijn1951}
N.~G. de Bruijn, \emph{On the number of positive integers $\le x$ and
free of prime factors $>y$}, Nederl. Akad. Wetensch. Proc. Ser. A
\textbf{54} (1951), 50--60.
\newline\url{https://pure.tue.nl/ws/files/1793853/597499.pdf}.

\bibitem{Erdos1976Manitoba}
P.~Erd\H{o}s, \emph{Problems and results on number theoretic properties of
consecutive integers and related questions}, Proceedings of the Fifth
Manitoba Conference on Numerical Mathematics (Winnipeg, 1975), Congressus
Numerantium XVI, Utilitas Math., Winnipeg, 1976, pp.~25--44.
\newline\url{https://users.renyi.hu/~p_erdos/1976-39.pdf}.

\bibitem{Erdos1976Debrecen}
P.~Erd\H{o}s, \emph{Problems and results on consecutive integers}, Publ.
Math. Debrecen \textbf{23} (1976), 271--282.
\newline\url{https://doi.org/10.5486/PMD.1976.23.3-4.15}.

\bibitem{ErdosPomerance1978}
P.~Erd\H{o}s and C.~Pomerance, \emph{On the largest prime factors of $n$
and $n+1$}, Aequationes Math. \textbf{17} (1978), 311--321.
\newline\url{https://users.renyi.hu/~p_erdos/1978-29.pdf}.

\bibitem{OpenAI2026}
OpenAI, \emph{The joint Dickman law for consecutive integers}, OpenAI
Math Release preprint, September~24, 2026, commit
\texttt{fd4aeeb2ee4f}.
\newline\url{https://github.com/openai/math/blob/fd4aeeb2ee4fc729c18d98444fed42fd0529eeeb/preprints/The-joint-Dickman-law-for-consecutive-integers-September-24-2026/paper.pdf}.

\bibitem{Schinzel1967}
A.~Schinzel, \emph{On two theorems of Gelfond and some of their
applications}, Acta Arith. \textbf{13} (1967), 177--236.
\newline\url{https://doi.org/10.4064/aa-13-2-177-236}.
\end{thebibliography}

\end{document}
